\chapter*{Notation and Conventions}
\addcontentsline{toc}{chapter}{Notation and Conventions}

You can skip this chapter and come back when a symbol surprises you. It is here
so that when Chapter~\ref{ch:gates} writes $B_t - R_t$ without comment, you know
what both terms mean and why the subtraction matters.

\section*{Time and Execution}

Time is discrete, indexed $t = 1, 2, \ldots$. An agent runs in steps, and at
each step it observes, decides, and possibly acts. Wall-clock time appears
separately when latency is the subject; until then, a ``step'' is one turn
through that loop.

\section*{State, Observation, and History}

\begin{itemize}
\item $S_t$ denotes the environment state at time $t$, which the agent generally cannot see.
\item $o_t$ denotes what the agent actually observes at time $t$.
\item $H_t = (o_1, a_1, \ldots, o_t)$ denotes the execution history through time $t$.
\end{itemize}

The distinction between $S_t$ and $o_t$ does real work in this book. Agents act
on $H_t$, a lossy and often stale record, while consequences land on $S_t$.

\section*{Actions and Policies}

\begin{itemize}
\item $a_t \in \mathcal{A}$ denotes the action taken at time $t$.
\item $\pi(\cdot \mid H_t)$ denotes the policy, mapping histories to actions.
\end{itemize}

Policies may be deterministic or stochastic, and may carry internal state of
their own (a budget tracker, a memory index, a set of open commitments).

\section*{Cost and Budgets}

\begin{itemize}
\item $c(e_t)$ denotes cost of event $e_t$.
\item $C(\tau)$ denotes total cost accumulated over a trace $\tau$.
\item $B_t$ denotes budget remaining at time $t$.
\item $R_t$ denotes budget \emph{reserved} for work the agent has not done yet but must.
\item $\hat{R}_t$ denotes the \emph{estimated} cost of completing the task from the
current state. It is an estimate, so it carries a distribution rather than a single number.
\item $\Delta_t = B_t - \hat{R}_t$ denotes \emph{slack}, the budget available after
setting aside what finishing is expected to cost.
\end{itemize}

Cost covers tokens, latency, external resource consumption, and the expected
cost of failure. Budgets are hard constraints unless a passage says otherwise.
$R_t$ is the term that separates an agent which finishes from one that runs out
of money three steps from the answer.

\section*{Traces}

A trace $\tau = \langle e_1, \ldots, e_T \rangle$ is the sequence of execution
events, namely observations, decisions, actions, gate outcomes, commits, and aborts.

Correctness, cost, and alignment are properties of traces. A single output can
look impeccable in a trace that violated an invariant four steps earlier, which
is why Chapter~\ref{ch:traces} treats the trace as the unit of evaluation.

\section*{Gates and Control Decisions}

A \emph{gate} is a control decision applied before commitment. Its outcomes are
\[
\{\textsc{Execute}, \textsc{Verify}, \textsc{Abstain}, \textsc{Refuse}\}.
\]

Gates read risk estimates, budgets, policies, and evidence. Chapter~\ref{ch:gates}
derives the thresholds rather than asserting them.

\section*{Speculation and Commitment}

\begin{itemize}
\item \emph{Speculation} denotes tentative execution whose results may be thrown away.
\item \emph{Commitment} denotes the moment results become irreversible and externally visible.
\end{itemize}

Speculative work is assumed abortable. Chapter~\ref{ch:speculation} spends
considerable effort on how often that assumption is false in practice, including
the case where an aborted branch still influences the model through retained
cache state.

\section*{Invariants and Policies}

\begin{itemize}
\item An \emph{invariant} must hold at every step of execution.
\item A \emph{policy} is a predicate constraining which actions are admissible.
\end{itemize}

Breaking an invariant counts as system failure even when the task succeeded.
A system that refunds the right customer by writing directly to the ledger has
still failed.

\section*{Conventions}

\begin{itemize}
\item Boldface marks defined terms on first use, and little else.
\item Probabilities and expectations condition on the agent's available
information unless stated otherwise.
\item Dollar figures and per-token prices are illustrative. Absolute prices move
fast; the \emph{ratios} between model tiers have been stable for years, and the
arguments in this book depend only on the ratios.
\item Examples are illustrative rather than exhaustive.
\end{itemize}

\tableofcontents

% ============================================================
% Part I: Foundations
% ============================================================
\mainmatter
